%! Composition of Functions — Unit 1, Lesson 1.5 — Group Activity (Answer Key)
\documentclass[10pt]{article}
\usepackage{precalculus-article}
\usepackage{precalculus-key}   % pulls in -boxes; do NOT also load -boxes
\newcommand{\TallMath}[1]{$\displaystyle #1\rule[-1.4em]{0pt}{3.2em}$}

\begin{document}

\pageheader{Unit 1, Lesson 1.5}{Group Activity --- Answer Key}

\vspace{0.12in}

\begin{scenariobox}[The Checkout Counter]{plumlight}
\small
A store is running a \$10-off coupon, and the state charges 6\% sales tax.
Both are functions of a dollar amount:

\smallskip
\begin{center}
\textbf{Coupon:} $c(p) = p - 10$ \qquad\qquad
\textbf{Sales tax:} $t(x) = 1.06x$
\end{center}

\smallskip
The register has to run \textit{both} --- but a register can only do one thing
at a time, so somebody had to choose an order. Every tier below uses these
same two functions. Round money answers to the nearest cent.
\end{scenariobox}

\vspace{0.1in}

\boxguard[30]
\begin{tcolorbox}[colback=white, colframe=black!40,
    title=\textbf{Tier R --- Remediate},
    fonttitle=\bfseries, arc=2mm, left=3mm, right=3mm, top=2mm, bottom=2mm]

The store's register applies the \textbf{coupon first}, then the tax. That is
the composite $t(c(p))$ --- inner function $c$, outer function $t$.

\begin{enumerate}[itemsep=8pt]
  \item \textbf{One machine at a time.} Fill the middle row first, all the way
        across, before you touch the bottom row.

        \smallskip
        \renewcommand{\arraystretch}{1.6}
        \begin{center}
        \begin{tabular}{|l|c|c|c|c|}
        \hline
        Sticker price $p$ & \$30 & \$50 & \$60 & \$80 \\
        \hline
        After the coupon: $c(p)$ & \ans{20} & \ans{40}
          & \ans{50} & \ans{70} \\
        \hline
        At the register: $t(c(p))$ & \ans{21.20} & \ans{42.40}
          & \ans{53.00} & \ans{74.20} \\
        \hline
        \end{tabular}
        \end{center}

  \item \textbf{Which machine ran first?} \quad
        \ans{$c$} \; / \; $t$ \; (circle one)

        In $t(c(p))$, the function written \textbf{last} runs
        \ans{first}.

  \item \textbf{Say what the middle row is.} The numbers you wrote in the
        middle row are not what anyone pays. In one sentence, what do they
        represent?

        \ansline{The price after the coupon comes off but before tax is added --- the number the tax machine gets.}

  \item A \$45 shirt goes through the same register. Work it in two steps.

        \begin{work}
          t(c(45)) &= t(45 - 10) \\
                   &= 1.06(35) \\
                   &= 37.10
        \end{work}

        The customer pays \$\ans{37.10}.
\end{enumerate}
\end{tcolorbox}

\vspace{0.1in}

\boxguard[30]
\begin{tcolorbox}[colback=white, colframe=black!40,
    title=\textbf{Tier A --- Approaching Proficiency},
    fonttitle=\bfseries, arc=2mm, left=3mm, right=3mm, top=2mm, bottom=2mm]

A second store rings up the \textbf{tax first}, then takes the coupon off ---
the composite $c(t(p))$. A \$60 jacket is on sale at both stores.

\begin{enumerate}[itemsep=8pt]
  \item Compute both totals, one step at a time.

        \begin{work}
          t(c(60)) &= t(60 - 10) \\
                   &= 1.06(50) \\
                   &= 53.00 \\
          c(t(60)) &= c(1.06 \cdot 60) \\
                   &= 63.60 - 10 \\
                   &= 53.60
        \end{work}

        Coupon first: \$\ans{53.00} \qquad Tax first: \$\ans{53.60}

  \item Build both composites as formulas in $p$.

        \begin{work}
          t(c(p)) &= 1.06(p - 10) \\
                  &= 1.06p - 10.60 \\
          c(t(p)) &= 1.06p - 10
        \end{work}

        $t(c(p)) = $ \ans{$1.06p - 10.60$} \qquad $c(t(p)) = $ \ans{$1.06p - 10$}

  \item Subtract the two formulas. The gap is \$\ans{0.60}, and it is the
        same at \textbf{every} price. Which store should the shopper use?

        \ansline{The first store --- coupon first is always \$0.60 cheaper, no matter the price.}

  \item Explain where that gap comes from --- what exactly is the second store
        charging tax on that the first store is not?

        \ansline{The second store taxes the full sticker price, including the \$10 the coupon takes off.}
        \ansline{Tax on that \$10 is $0.06(10) = \$0.60$ --- exactly the gap.}
\end{enumerate}
\end{tcolorbox}

\vspace{0.1in}

\boxguard[30]
\begin{tcolorbox}[colback=white, colframe=black!40,
    title=\textbf{Tier E --- Extension},
    fonttitle=\bfseries, arc=2mm, left=3mm, right=3mm, top=2mm, bottom=2mm]

\begin{enumerate}[itemsep=8pt]
  \item \textbf{Swap the coupon for a percentage.} The store replaces the \$10
        coupon with 25\% off: $d(p) = 0.75p$. A \$40 item goes through.

        \begin{work}
          d(t(40)) &= d(42.40) \\
                   &= 0.75(42.40) \\
                   &= 31.80 \\
          t(d(40)) &= t(30) \\
                   &= 1.06(30) \\
                   &= 31.80
        \end{work}

        $d(t(40)) = $ \$\ans{31.80} \qquad $t(d(40)) = $ \$\ans{31.80}

        These two \textbf{are} equal. So ``order matters'' is a warning, not a
        law --- some pairs commute. In one sentence, what is different about
        $d$ and $t$ that lets them trade places, when $c$ and $t$ cannot?

        \ansline{$d$ and $t$ both just multiply, and multiplication can be done in any order: $0.75(1.06p) = 1.06(0.75p)$.}
        \ansline{The coupon $c$ subtracts, and multiplying after subtracting is not the same as subtracting after multiplying.}

  \item \textbf{Read the receipt backwards.} A register prints the rule
        $R(p) = 1.06p - 10.60$. Which composite is it?

        \smallskip
        \ans{$R = t \circ c$} \; / \; $R = c \circ t$ \quad (circle one)

        \smallskip
        Justify by rebuilding it: $R(p) = $ \ans{$1.06(p - 10) = 1.06p - 10.60$}

  \item \textbf{The do-nothing function.} A register in a tax-free state runs
        $I(p) = p$ before the coupon.

        \smallskip
        $c(I(p)) = $ \ans{$p - 10$} \qquad $I(c(p)) = $ \ans{$p - 10$}

        \smallskip
        A function that leaves every other function alone is called the
        \ans{identity} function.

  \item \textbf{Back to Lesson 1.4.} You wrote $g(x) = f(x + 3)$ last lesson
        and called it a horizontal shift. Write it as a composition:

        \smallskip
        $g(x) = f(h(x))$ \quad where \quad $h(x) = $ \ans{$x + 3$}

        \smallskip
        So a change made \textit{inside} $f$ is really a
        \ans{composition} --- which is why it acts on the inputs.
\end{enumerate}
\end{tcolorbox}

\end{document}
